\clearpage
\SourceLeaf{412}
\begin{center}\large ZUR TRANSFORMATION DER FLÄCHEN.\end{center}
\noindent\hspace*{1.2em}Die Punkte einer an sich gegebenen Fläche, deren Form aber noch unbestimmt bleibt, werden durch $r$ und $\theta$ bestimmt, wo $r$ die Entfernung von einem gegebenen Punkt, gemessen in kürzester Linie, und $\theta$ den Winkel dieser kürzesten Linie mit einer constanten andern, gleichfalls durch jenen Punkt gehenden [kürzesten Linie] bedeutet. Jede Linie, für die $\theta$ constant ist, wird dann von jeder andern, für die $r$ constant ist, unter einem rechten Winkel geschnitten. Setzt man ein Element der letztern $=md\theta$, so ist $m$ bloss von der Natur der Fläche, aber nicht von ihrer Form abhängig, und also eine gegebene Function von $r$ und $\theta$.

Sobald man der Fläche eine bestimmte Form gegeben, findet folgendes statt für eine Linie, für die $\theta$ constant ist.

Es seien $x,y,z$ indefinite Coordinaten eines Punkts derselben, ferner
\[
\begin{array}{lll}
\displaystyle\frac{\partial x}{\partial r}=\xi,&\displaystyle\frac{\partial y}{\partial r}=\eta,&\displaystyle\frac{\partial z}{\partial r}=\zeta,\\[6pt]
\displaystyle\frac{\partial\xi}{\partial r}=pX,&\displaystyle\frac{\partial\eta}{\partial r}=pY,&\displaystyle\frac{\partial\zeta}{\partial r}=pZ,
\end{array}
\]
\noindent wo
\[
p=\BareRoot{\left(\frac{\partial\xi^2}{\partial r^2}+\frac{\partial\eta^2}{\partial r^2}+\frac{\partial\zeta^2}{\partial r^2}\right)},
\]
\noindent also $\dfrac1p$ Krümmungshalbmesser und
\[
XX+YY+ZZ=1.
\]
